\section{The Existing Holm-Bonferroni Implementation}
\label{sec:impl}

\subsection{Current Implementation in bisect-analysis}

The \bisect{} codebase already implements Holm-Bonferroni correction
in the bloc voting analysis module.
The relevant code is in:
\begin{verbatim}
crates/bisect-analysis/src/bloc_voting.rs
\end{verbatim}

The \texttt{BlocVotingResult} struct includes:
\begin{verbatim}
pub struct BlocVotingResult {
    pub district_id: u32,
    pub metric: String,
    pub p_value_raw: f64,
    pub p_value_bonferroni: f64,
    pub p_value_holm: f64,  // Holm-Bonferroni corrected
    pub significant_holm: bool,
}
\end{verbatim}

The \texttt{p\_value\_holm} field is computed using the standard
Holm-Bonferroni algorithm: sort p-values in ascending order, then
for each p-value $p_{(k)}$, the Holm-corrected p-value is:
\[
  p_{(k)}^{\rm Holm} = \max_{j \leq k} \min\!\left[(m - j + 1) p_{(j)}, 1\right].
\]

\subsection{Why Holm Is Appropriate for Bloc Voting but Not L-Track}

The Holm-Bonferroni procedure controls the FWER, which is appropriate
for bloc voting analysis because:
\begin{enumerate}
\item The number of tests per state is small (typically $m = 5$--$20$
  districts tested for VRA compliance).
\item A false positive in VRA analysis has serious legal consequences:
  incorrectly certifying a district as a majority-minority district when
  it is not could constitute a VRA violation.
\item FWER control ensures that the probability of any false VRA claim
  is $\leq \alpha$.
\end{enumerate}

For the L-track portfolio analysis, these conditions do not hold:
\begin{enumerate}
\item The number of tests is large ($m = 1{,}050$), making FWER control
  overly conservative.
\item The consequence of a false positive (one borderline state is
  flagged as potentially gerrymandered when it is not) is manageable:
  further litigation will resolve the false positive.
\item Many hypotheses are expected to be true (the 2010 and 2020 cycles
  include genuinely gerrymandered maps), making FDR control more
  powerful.
\end{enumerate}

\subsection{S.4's Extension}

S.4 extends the existing implementation by adding a BH FDR module to
\texttt{bisect-analysis}:
\begin{verbatim}
crates/bisect-analysis/src/fdr_correction.rs
\end{verbatim}

The module implements:
\begin{itemize}
\item Benjamini-Hochberg FDR at configurable level $q$
\item Benjamini-Yekutieli correction for dependent tests
\item Q-value computation~\citep{storey2002direct} for adaptive FDR estimation
\item Integration with the L-track metric battery
\end{itemize}

The module is invoked via:
\begin{verbatim}
bisect analyze --types partisan-fdr --year 2020
\end{verbatim}
